%% annex/p03/qe-categoricity-games.tex --- P3 QE and back-and-forth proofs.
%% Individual items retain roughly 40--60% of their v1.1 proof-line counts.

\begin{annexitem}{annex:p3.qe.1}{thm:p03-model-completeness-reduction}{Reduction of model completeness to inclusions}
\begin{proof}
Given an embedding \(h:\mathcal M\to\mathcal N\), corestrict it to the isomorphism \(h_0:\mathcal M\to h[M]\).
Its image is a substructure of \(\mathcal N\) and, by isomorphism invariance, a model of \(T\); hence \(h[M]\preccurlyeq\mathcal N\).
For every formula \(\varphi\) and tuple \(\bar a\in M\),
\[
 \mathcal M\entails\varphi(\bar a)\ \Longleftrightarrow\ h[M]\entails\varphi(h\bar a)\ \Longleftrightarrow\ \mathcal N\entails\varphi(h\bar a).
\]
Thus the arbitrary embedding \(h\) is elementary.
\end{proof}
\end{annexitem}

\begin{annexitem}{annex:p3.qe.2}{thm:p03-dlo-complete-model-complete}{Completeness and model completeness of DLO}
\sketchnote{omitted: DNF and existential distribution; equality substitution
and conflicts; bound compatibility; finite disequality avoidance with both,
lower-only, upper-only, or no bounds; and the Boolean induction.}
QE sends every DLO sentence to a quantifier-free sentence.  Since \(\{<\}\) has no closed terms or atomic sentences, its truth value is fixed; first-order adequacy gives deductive completeness.
If DLO models satisfy \(\mathcal M\subseteq\mathcal N\), quantifier-free absoluteness and QE give \(\mathcal M\preccurlyeq\mathcal N\).
Theorem~\ref{thm:p03-model-completeness-reduction} then covers arbitrary embeddings.
\end{annexitem}

\begin{annexitem}{annex:p3.qe.3}{thm:p03-back-and-forth-test}{Back-and-forth test}
\begin{proof}
Enumerate \(M=\{a_n:n\in\N\}\) and \(N=\{b_n:n\in\N\}\), starting at the prescribed \(f_0\).
At stage \(2n\) apply forth to \(a_n\); at stage \(2n+1\) apply back to \(b_n\), yielding a chain in \(\mathcal F\) above \(f_0\).
Its union \(h\) is injective; alternation makes it total and onto.
Each atomic instance occurs at a finite stage, so \(h\) preserves and reflects it.
Thus \(h\) is an isomorphism extending the arbitrary \(f_0\).
\end{proof}
\end{annexitem}

\begin{annexitem}{annex:p3.qe.4}{thm:p03-countable-saturated-automorphism}{Finite elementary self-maps of countable saturated structures}
\begin{proof}
Let \(\mathcal F\) be the finite elementary self-maps of \(\mathcal M\) extending \(f\); it is nonempty, and its members are partial isomorphisms.
Theorem~\ref{thm:p03-saturation-homogeneity} extends any \(g\in\mathcal F\) to include a prescribed domain element.
For a prescribed range element, apply that theorem to \(g^{-1}\) and invert the extension.
Both operations still extend \(f\), so \(\mathcal F\) satisfies forth and back.
Theorem~\ref{thm:p03-back-and-forth-test}, started at \(f\), yields the required automorphism.
\end{proof}
\end{annexitem}

\begin{annexitem}{annex:p3.qe.5}{thm:p03-dlo-countable-categorical}{Countable categoricity of dense linear order}
\begin{proof}
Let \(\mathcal F\) be the finite order-preserving partial bijections between the DLO models; it contains the empty map.
Covered elements need no change.  A new element's cut has four forms: two bounds, lower only, upper only, or none.
Choose a fresh image, respectively, by density, no right endpoint, no left endpoint, or nonemptiness.
This gives forth extensions; the symmetric four cases give back extensions.
Theorem~\ref{thm:p03-back-and-forth-test} yields an isomorphism.
\end{proof}
\end{annexitem}
